\documentclass[11pt,compress,t,notes=noshow, xcolor=table]{beamer}

\input{../../style/preamble}
\input{../../latex-math/basic-math}
\input{../../latex-math/basic-ml}
\input{../../latex-math/optim-basics}

\title{Optimization in Machine Learning}

\begin{document}

\titlemeta{
Constrained Optimization
}{
Introduction
}{
figure/intro_feasible_region.pdf
}{
\item Optimization over a feasible region
\item Why unconstrained updates can leave feasibility
\item Primal and dual viewpoints
\item Constraint structure: convex / linear versus norm constraints
}

\begin{framev}[fs=footnotesize]{Optimization Over a Feasible Region}
\textbf{General form}
$$
\min_{\xv \in \R^d} f(\xv)
\qquad
\text{s.t. } \qquad
g_i(\xv) \le 0,\; i=1,\dots,m,
\qquad
h_j(\xv) = 0,\; j=1,\dots,\ell.
$$
\textbf{Feasible set}
$$
\S := \left\{\xv \in \R^d \mid g_i(\xv) \le 0,\; h_j(\xv)=0 \right\}.
$$
\spacer
\splitV[0.50]{
\begin{itemize}
\item Constrained optimization minimizes $f$ only over $\S$, not over all of $\R^d$.
\item The unconstrained minimizer can lie outside $\S$, so the constrained optimum often lies on the boundary.
\item Toy problem for this chapter:
\begin{align*}
\min_{\wv \in [0,1]\times[0,0.9]} (w_1 - 1.2)^2 &+ (w_1 - 1.2)(w_2 - 0.2)\\
&+ (w_2 - 0.2)^2.
\end{align*}
\end{itemize}
}{
\image[1]{figure/intro_feasible_region.pdf}
}
\end{framev}


\begin{framev}[fs=footnotesize]{Why Unconstrained Methods Need Adaptation}
\textbf{Vanilla gradient step}
$$
\zv^{[t+1]} = \xv^{[t]} - \alpha_t \nabla f(\xv^{[t]}).
$$
\vfill
\splitV[0.50]{
\begin{itemize}
\item Even if $\xv^{[t]} \in \S$, the tentative point $\zv^{[t+1]}$ may satisfy $\zv^{[t+1]} \notin \S$.
\item The same issue appears for coordinate descent and Newton-type methods: their natural unconstrained updates need not respect the constraints.
\item Constrained optimization therefore modifies the update rule, repairs an infeasible trial point, or reformulates the problem.
\end{itemize}
}{
\imageC[0.96]{figure/intro_unconstrained_step.pdf}
}
\end{framev}


\begin{framev}[fs=footnotesize]{Two Viewpoints}
\textbf{Primal methods}
\begin{itemize}
\item work directly in the original variables
\item keep iterates feasible or repair a tentative step
\end{itemize}
\textbf{Dual methods}
\begin{itemize}
\item relax constraints with Lagrange multipliers
\item optimize the resulting problem in multiplier space
\end{itemize}
\textbf{Bridge methods}
Penalty and barrier methods sit between both viewpoints and connect directly to primal-dual interior-point ideas.
\imageC[0.60]{figure/intro_primal_dual_roadmap.pdf}
\end{framev}


\begin{framev}[fs=footnotesize]{Constraint Structure Matters}
\textbf{There is no single best constrained optimizer.}
\begin{itemize}
\item The right method depends strongly on how the feasible set is described and how expensive it is to enforce.
\item \textbf{Linear / convex constraints:} affine sets, boxes, simplices, and general convex regions often admit projections, coordinate slices, and duality / KKT structure.
\item \textbf{Norm constraints:} balls, spheres, and orthogonality constraints lead to special geometry used in eigenvalue problems and PCA.
\end{itemize}
\imageC[0.98]{figure/intro_constraint_classes.pdf}
\end{framev}

\endlecture
\end{document}
